\subsection{Posición de las Llaves}

\subsubsection{Regla: estilo Allman}

\textbf{Regla:} La llave de apertura y la llave de cierre de cualquier bloque (clases, métodos, \texttt{if}, \texttt{else}, \texttt{for}, \texttt{switch}, entre otros) deberán colocarse cada una en su propia línea, alineadas con el nivel de indentación del bloque que delimitan. Queda prohibido que un \texttt{else}, \texttt{catch} o \texttt{finally} comparta línea con la llave de cierre del bloque anterior.

\textbf{Descripción:} Las convenciones de código de C\# de Microsoft especifican el estilo ``Allman'' para las llaves, en el que cada llave de apertura y cierre ocupa su propia línea (Microsoft, 2025e).

\textbf{Código correcto:}
\begin{lstlisting}[style=csharpstyle]
if (player.IsActive)
{
    ExecuteTurn(player);
}
else
{
    SkipTurn(player);
}
\end{lstlisting}

\textbf{Código incorrecto:}
\begin{lstlisting}[style=csharpstyle]
if (player.IsActive) {
    ExecuteTurn(player);
} else {
    SkipTurn(player);
}
\end{lstlisting}
